% SPDX-License-Identifier: Apache-2.0
\section{The Lowering, Checked}
\label{sec:x-checked}

A lowering that is only printed is a claim. From this section on, a
lowered unit is also run, and against the same stimulus the Rust
simulation had. The build does it: the example is run with
\code{TXHDL\_VHDL} set, so beside its trace it writes the unit's VHDL,
which \code{\#[lower]} renders from the same statements as the
Verilog, and a ports file; \code{//tools/fst2tb} reads the trace and
the ports and writes a testbench that replays the trace's inputs into
the entity and asserts, tick by tick, that the entity's registers and
outputs are what the trace holds; and \code{vhdl\_test} from
\code{rules\_nvc} simulates it under \code{nvc}, built from source,
hermetically. The test passes only if every comparison held.

The testbench's timing is the model of time, made literal. One tick is
one nanosecond and the clock rises at even ticks. An input the trace
holds at tick $2k$ was set before that edge, so it is applied at
$2k-1$. A register the trace holds at $2k$ took its value at that
edge, so it is checked at $2k+1$, through a VHDL-2008 external name.
A wire the Rust process drove at $2k$ was computed from the registers
as they stood before that edge, so it is checked at $2k-1$, where the
entity's continuous assignment shows the same value. That last rule
is a fact about the runtime worth knowing: a wire a process drives
inside its iteration lags the hardware's continuous assignment by the
process's own step.

The counter of the cheat sheet, the blinky, the sequencer and the ALU
pass.
A negative control was run once: the counter's adder changed to add
two makes the test fail at the first tick, so a pass means something.

The stage of Section~\ref{sec:x-stage} is not checked, and the reason
is a finding. The runtime's channel is a mailbox: an offer sits until
taken, and the consumer takes it in the same step it is offered, so
\code{valid} is never high at a tick's end and \code{ready} means the
slot is empty. The lowering's channel is a valid/ready handshake held
across edges. The two describe the same design but not the same
signals per tick, so a trace of one is not a check of the other. To
check a unit with channels, the runtime's channel has to become a
handshake that is cycle-accurate, valid held until ready is seen at an
edge; that is the next change to the runtime, and it will change the
FIFO's and the arbiter's traces.
